\section{Introduction}

Molecular property prediction is a central component of AI-driven drug discovery (AIDD), supporting early-stage tasks such as ADMET estimation, toxicity screening, blood--brain barrier permeability prediction, and lead prioritization. Public benchmarks have accelerated progress by standardizing datasets, molecular representations, and evaluation metrics. MoleculeNet established a widely used benchmark suite for molecular machine learning \citep{wu2018moleculenet}, while Therapeutics Data Commons (TDC) broadened therapeutics machine learning by organizing datasets, tasks, splits, and evaluation protocols across drug discovery and development \citep{huang2021tdc}. In parallel, learned molecular representations and graph-based models have shown strong performance in molecular property prediction, but their interpretation still depends on how generalization is defined and evaluated \citep{feinberg2018potentialnet,yang2019learned}.

A benchmark score is not only a property of the model. It is also a property of the split. When molecules are randomly divided into training and test sets, chemically similar compounds can appear on both sides of the split, potentially making test performance appear stronger than it would be under more distant chemical generalization. Scaffold splitting is commonly used to address this issue by grouping molecules according to molecular frameworks related to Bemis--Murcko scaffolds \citep{bemis1996properties}. In principle, scaffold splitting asks a more chemically meaningful question: whether a model can generalize from known scaffold families to held-out scaffold families.

Scaffold splitting is nevertheless not a complete solution to benchmark reliability. A scaffold split can remove direct scaffold overlap while simultaneously changing the target distribution, concentrating the test set into a small number of scaffold groups, or leaving chemically similar molecules across apparently different scaffolds. Recent work has argued that scaffold splits can still overestimate virtual screening performance because different scaffolds may remain chemically similar \citep{guo2024scaffold}. These concerns are closely related to long-standing QSAR validation and applicability-domain principles, where external validation, chemical coverage, and domain boundaries are central to model interpretation \citep{tropsha2003earnest,gramatica2007principles,netzeva2005applicability,sahigara2012applicability}.

This study treats the train/test split as an object of measurement. Random, ordinary scaffold, and target-balanced scaffold protocols are compared across six public molecular property datasets. The target-balanced scaffold split is used as a diagnostic counterfactual: it preserves scaffold-level separation while attempting to reduce the train/test target mean gap. It is not proposed as a universal replacement for scaffold splitting. Instead, it helps test whether an observed scaffold gap remains after one major confounder, target shift, is reduced.

The contribution is a reproducible split-diagnostic audit rather than a new molecular predictor. Model performance gaps, target distribution shifts, scaffold statistics, test-to-train Tanimoto similarity, outlier cases, and model-ranking changes are reported jointly. This combination of diagnostics provides a more transparent interpretation of benchmark results: it reveals when random splits are chemically close, when scaffold splits are statistically shifted, when balancing helps, and when model rankings depend on the split protocol. The central claim is deliberately limited but practically important: in molecular property prediction benchmarks, split diagnostics should accompany model scores before conclusions about molecular generalization are drawn.
